% Chapter 1, Topic from _Linear Algebra_ Jim Hefferon
%  http://joshua.smcvt.edu/linalg.html
%  2001-Jun-09
\topic{计算机代数系统}
\index{computer algebra systems|(}
本章中的线性方程组规模都很小，
用手求解很容易。
对于大型方程组，包括含有成千上万个方程的方程组，
我们就需要计算机了。
为此有一些专用程序，如 LINPACK\index{LINPACK}。
流行的还有通用计算机代数系统，包括 \Maple，\index{Maple} 
\textit{Mathematica}，\index{Mathematica} 或 \textit{MATLAB}，\index{MATLAB} 
以及 \Sage\index{Sage}。

例如，在关于网络的专题中，我们需要求解下面这个方程组。
\begin{equation*}
  \begin{linsys}{7}
    i_0  &-  &i_1  &-  &i_2  &   &    &  &    &   &    &  &    &=  &0  \\
         &   &i_1  &   &     &-  &i_3 &  &    &-  &i_5 &  &    &=  &0  \\
         &   &     &   &i_2  &   &    &- &i_4 &+  &i_5 &  &    &=  &0  \\
         &   &     &   &     &   &i_3 &+ &i_4 &   &    &- &i_6 &=  &0  \\
         &   &5i_1 &   &     &+  &10i_3  &  & &   &    &  &    &=  &10  \\
         &   &     &   &2i_2 &   &    &+ &4i_4 &  &    &  &    &=  &10  \\
         &   &5i_1 &-  &2i_2 &   &    &  &    &+  &50i_5 &&    &=  &0   
  \end{linsys}
\end{equation*}
用手工来算既费时又容易出错。
用计算机更好。

下面是用 \textit{Sage} 求解该方程组的过程。
（做法有很多种；这里采用的方法胜在简单。）
\begin{lstlisting}
sage: var('i0,i1,i2,i3,i4,i5,i6')
(i0, i1, i2, i3, i4, i5, i6)
sage: network_system=[i0-i1-i2==0, i1-i3-i5==0, 
....:       i2-i4+i5==0, i3+i4-i6==0, 5*i1+10*i3==10,
....:       2*i2+4*i4==10, 5*i1-2*i2+50*i5==0]
sage: solve(network_system, i0,i1,i2,i3,i4,i5,i6)     
[[i0 == (7/3), i1 == (2/3), i2 == (5/3), i3 == (2/3), 
      i4 == (5/3), i5 == 0, i6 == (7/3)]] 
\end{lstlisting}
真是神奇。

下面是在 Maple 中求解同一个方程组。
我们输入系数数组
与常数向量，
然后就得到解。
\begin{lstlisting}
> A:=array( [[1,-1,-1,0,0,0,0],
             [0,1,0,-1,0,-1,0],
             [0,0,1,0,-1,1,0],
             [0,0,0,1,1,0,-1],
             [0,5,0,10,0,0,0],
             [0,0,2,0,4,0,0],
             [0,5,-2,0,0,50,0]] );
> u:=array( [0,0,0,0,10,10,0] );
> linsolve(A,u);
      7  2  5  2  5     7
    [ -, -, -, -, -, 0, - ]
      3  3  3  3  3     3
\end{lstlisting}

若方程组有无穷多解，
则程序会返回一个参数化。


\begin{exercises}
  \item 
    用计算机求解本章开头提出的
    两个问题。
    \begin{exparts}
      \partsitem 这是静力学问题。
         \begin{align*}
            40h+15c  &= 100  \\
            25c      &= 50+50h
         \end{align*}
      \partsitem 这是化学问题。
         \begin{align*} 
             7h      &= 7j  \\
             8h +1i  &= 5j+2k  \\
             1i      &= 3j  \\
             3i      &= 6j+1k
         \end{align*}
    \end{exparts}
    \begin{answer}
      \begin{exparts}
        \partsitem 
\Sage{} 求解如下。
\begin{lstlisting}
sage: var('h,c')
(h, c)
sage: statics = [40*h + 15*c == 100,
....:            25*c == 50 + 50*h]
sage: solve(statics, h,c)
[[h == 1, c == 4]] 
\end{lstlisting}
其他计算机代数系统也有类似的命令。
下面这些 \Maple{} 命令
\begin{lstlisting}
> A:=array( [[40,15],
             [-50,25]] );
> u:=array([100,50]);
> linsolve(A,u);
\end{lstlisting}
           给出答案 \lstinline![1,4]!。
        \partsitem 这里有一个自由变量。
\Sage{} 给出如下结果。
\begin{lstlisting}
sage: var('h,i,j,k')
(h, i, j, k)
sage: chemistry = [7*h == 7*j,
....:              8*h + 1*i == 5*j + 2*k,
....:              1*i == 3*j,
....:              3*i == 6*j + 1*k]
sage: solve(chemistry, h,i,j,k)
[[h == 1/3*r1, i == r1, j == 1/3*r1, k == r1]]          
\end{lstlisting}
类似地，下面这个 \Maple{} 会话
\begin{lstlisting}
> A:=array( [[7,0,-7,0],
             [8,1,-5,2],
             [0,1,-3,0],
             [0,3,-6,-1]] );
> u:=array([0,0,0,0]);
> linsolve(A,u);
\end{lstlisting}
         给出同样的回答（但参数为 $t_1$）。
      \end{exparts}
    \end{answer}
\item 
    用计算机求解这些来自
    第一小节的方程组，
    或者得出“无穷多解”或“无解”的结论。
    \begin{exparts*}
      \partsitem \(
               \begin{linsys}[t]{2}
                  2x  &+  &2y  &=  &5  \\
                   x  &-  &4y  &=  &0  
               \end{linsys}
             \)
      \partsitem \(
               \begin{linsys}[t]{2}
                  -x  &+  &y   &=  &1  \\
                   x  &+  &y   &=  &2  
               \end{linsys}
             \) 
      \partsitem  \(
               \begin{linsys}[t]{3}
                   x  &-  &3y  &+  &z  &=  &1  \\
                   x  &+  &y   &+  &2z &=  &14 
                \end{linsys}
             \) 
      \partsitem  \(
               \begin{linsys}[t]{2}
                  -x  &-  &y   &=  &1  \\
                 -3x  &-  &3y  &=  &2  
               \end{linsys}
             \) 
      \partsitem  \(
               \begin{linsys}[t]{3}
                       &   &4y  &+  &z  &=  &20 \\
                   2x  &-  &2y  &+  &z  &=  &0  \\
                    x  &   &    &+  &z  &=  &5  \\
                    x  &+  &y   &-  &z  &=  &10 
                \end{linsys}
             \)
      \partsitem \( \begin{linsys}[t]{4}
                 2x  &   &   &+  &z  &+  &w  &=  &5  \\
                     &   &y  &   &   &-  &w  &=  &-1 \\
                 3x  &   &   &-  &z  &-  &w  &=  &0  \\
                 4x  &+  &y  &+  &2z &+  &w  &=  &9  
               \end{linsys}
            \)
    \end{exparts*}
    \begin{answer}
      \begin{exparts}
        \partsitem 答案为 \( x=2 \) 与 \( y=1/2 \)。
          用 \Sage{} 会话求解如下。
\begin{lstlisting}
sage: var('x,y')
(x, y)
sage: system = [2*x + 2*y == 5,
....:              x - 4*y == 0]
sage: solve(system, x,y)
[[x == 2, y == (1/2)]]          
\end{lstlisting}
          下面这个 \Maple{} 会话
\begin{lstlisting}
> A:=array( [[2,2],
             [1,-4]] );
> u:=array([5,0]);
> linsolve(A,u);            
\end{lstlisting}
          当然得到相同的答案。
        \partsitem 答案为 \( x=1/2 \) 与 \( y=3/2 \)。 
\begin{lstlisting}
sage: var('x,y')
(x, y)
sage: system = [-1*x + y == 1,
....:              x + y == 2]
sage: solve(system, x,y)
[[x == (1/2), y == (3/2)]]          
\end{lstlisting}
        \partsitem 该方程组有无穷多解。
           在第一小节中，取 $z$ 为参数，
           我们得到 $x=(43-7z)/4$ 与 $y=(13-z)/4$。
           \Sage{} 得到相同的结果。
\begin{lstlisting}
sage: var('x,y,z')
(x, y, z)
sage: system = [x - 3*y + z == 1,
....:           x + y + 2*z == 14]
sage: solve(system, x,y)
[[x == -7/4*z + 43/4, y == -1/4*z + 13/4]]             
\end{lstlisting}
           \Maple{} 回答 $(-12+7t_1, t_1, 13-4t_1)$，
           它更愿意取 $y$ 作参数。
        \partsitem 该方程组无解。
           \Sage{} 给出空列表。
\begin{lstlisting}
sage: var('x,y')
(x, y)
sage: system = [-1*x - y == 1,
....:           -3*x - 3*y == 2]
sage: solve(system, x,y)
[]          
\end{lstlisting}
           类似地，
           当把数组 $A$ 与向量 $u$ 交给 \Maple{}
           并要求它执行 \lstinline!linsolve(A,u)! 时，
           它完全不返回任何结果；也就是说，它以
           无解作答。
        \partsitem \Sage{} 求得 
\begin{lstlisting}
sage: var('x,y,z')
(x, y, z)
sage: system = [       4*y + z == 20,
....:            2*x - 2*y + z == 0,
....:             x +        z == 5,
....:             x +    y - z == 10]
sage: solve(system, x,y,z)
[[x == 5, y == 5, z == 0]]          
\end{lstlisting}
           即解为 \( (x,y,z)=(5,5,0) \)。
        \partsitem 解有无穷多个。
           \Sage{} 求解如下。
\begin{lstlisting}
sage: var('x,y,z,w')
(x, y, z, w)
sage: system = [ 2*x +        z + w == 5,
....:                   y -       w == -1,
....:            3*x -        z - w == 0,
....:            4*x +  y + 2*z + w == 9]
sage: solve(system, x,y,z,w)
[[x == 1, y == r2 - 1, z == -r2 + 3, w == r2]]          
\end{lstlisting}
           \Maple{} 给出 $(1, -1+t_1, 3-t_1, t_1)$。
      \end{exparts}
    \end{answer}
  \item 
    用计算机求解第二小节中的这些方程组。
    \begin{exparts*}
      \partsitem \( \begin{linsys}[t]{2}
                  3x  &+  &6y  &=  &18  \\
                   x  &+  &2y  &=  &6   
                   \end{linsys}  \)
      \partsitem \( \begin{linsys}[t]{2}
                   x  &+  &y   &=  &1  \\
                   x  &-  &y   &=  &-1   
                    \end{linsys}  \)
      \partsitem \( \begin{linsys}[t]{3}
                   x_1  &   &     &+  &x_3   &=  &4  \\
                   x_1  &-  &x_2  &+  &2x_3  &=  &5  \\
                  4x_1  &-  &x_2  &+  &5x_3  &=  &17  
                   \end{linsys}  \)
      \partsitem \( \begin{linsys}[t]{3}
                   2a   &+  &b    &-  &c     &=  &2  \\
                   2a   &   &     &+  &c     &=  &3  \\
                    a   &-  &b    &   &      &=  &0   
                    \end{linsys}  \)
      \partsitem \( \begin{linsys}[t]{4}
                     x  &+  &2y   &-   &z   &    &    &=  &3  \\
                    2x  &+  &y    &    &    &+   &w   &=  &4  \\
                     x  &-  &y    &+   &z   &+   &w   &=  &1  
                    \end{linsys}  \)
      \partsitem \( \begin{linsys}[t]{4}
                     x  &   &     &+   &z   &+   &w   &=  &4  \\
                    2x  &+  &y    &    &    &-   &w   &=  &2  \\
                    3x  &+  &y    &+   &z   &    &    &=  &7  
                     \end{linsys}  \)
    \end{exparts*}
    \begin{answer}
      \begin{exparts}
        \partsitem 该方程组有无穷多解。
              在第二小节中我们把解集写成
              \begin{equation*}
              \set{\colvec{6 \\ 0}+\colvec{-2 \\ 1}y
                      \suchthat y\in\Re}
              \end{equation*}
              \Sage{} 求得 
\begin{lstlisting}
sage: var('x,y')
(x, y)
sage: system = [ 3*x + 6*y == 18,
....:              x + 2*y == 6]
sage: solve(system, x,y)
[[x == -2*r3 + 6, y == r3]]                
\end{lstlisting}
              而 \Maple{} 回答 $(6-2t_1,t_1)$。
        \partsitem 解集只有一个元素。
          \begin{equation*}
             \set{\colvec{0 \\ 1} }
          \end{equation*}
          \Sage{} 给出如下结果。
\begin{lstlisting}
sage: var('x,y')
(x, y)
sage: system = [ x + y == 1,
....:            x - y == -1]
sage: solve(system, x,y)
[[x == 0, y == 1]]            
\end{lstlisting}
        \partsitem 该方程组的解集是无限的
          \begin{equation*}
            \set{\colvec{4 \\ -1 \\ 0}+\colvec{-1 \\ 1 \\ 1}x_3
                             \suchthat x_3\in\Re}
          \end{equation*}
          \Sage{} 给出如下结果
\begin{lstlisting}
sage: var('x1,x2,x3')
(x1, x2, x3)
sage: system = [   x1 +        x3 == 4,
....:              x1 - x2 + 2*x3 == 5,
....:            4*x1 - x2 + 5*x3 == 17]
sage: solve(system, x1,x2,x3)
[[x1 == -r4 + 4, x2 == r4 - 1, x3 == r4]]            
\end{lstlisting}
          而 Maple 给出 $(t_1,-t_1+3,-t_1+4)$。
        \partsitem 存在唯一解
           \begin{equation*}
             \set{\colvec{1 \\ 1 \\ 1}}
           \end{equation*}
           \Sage{} 也求出了它。
\begin{lstlisting}
sage: var('a,b,c')
(a, b, c)
sage: system = [ 2*a +  b - c == 2,
....:            2*a +      c == 3,
....:              a - b      == 0]
sage: solve(system, a,b,c)
[[a == 1, b == 1, c == 1]]             
\end{lstlisting}
        \partsitem 该方程组有无穷多解；在
           第二小节中我们用两个参数描述了
           该解集。
           \begin{equation*}
             \set{\colvec{5/3 \\ 2/3 \\ 0 \\ 0}
                  +\colvec{-1/3 \\ 2/3 \\ 1 \\ 0}z
                  +\colvec{-2/3 \\ 1/3 \\ 0 \\ 1}w
                  \suchthat z,w\in\Re}
           \end{equation*}
           \Sage{} 也给出同样的结果。
\begin{lstlisting}
sage: var('x,y,z,w')
(x, y, z, w)
sage: system = [   x + 2*y - z     == 3,
....:            2*x +   y     + w == 4,
....:              x -   y + z + w == 1]
sage: solve(system, x,y,z,w)
[[x == r6, y == -r5 - 2*r6 + 4, z == -2*r5 - 3*r6 + 5, w == r5]]             
\end{lstlisting}
           \Maple{} 也是如此：$(3-2t_1+t_2,t_1,t_2,-2+3t_1-2t_2)$。
        \partsitem 解集为空。
\begin{lstlisting}
sage: var('x,y,z,w')
(x, y, z, w)
sage: system = [   x +      z + w == 4,
....:            2*x +  y     - w == 2,
....:            3*x +  y + z     == 7]
sage: solve(system, x,y,z,w)
[]          
\end{lstlisting}
      \end{exparts}
    \end{answer}
  \item 
    对于一般的
    $\nbyn{2}$ 方程组，计算机给出的解是什么？
    \begin{equation*}
      \begin{linsys}{2}
        ax  &+  &cy  &=  &p  \\
        bx  &+  &dy  &=  &q
      \end{linsys}
    \end{equation*}
    \begin{answer}
      \Sage{} 求解如下。
\begin{lstlisting}
sage: var('x,y,a,b,c,d,p,q')
(x, y, a, b, c, d, p, q)
sage: system = [ a*x +  c*y == p,
....:            b*x +  d*y == q]
sage: solve(system, x,y)
[[x == -(d*p - c*q)/(b*c - a*d), y == (b*p - a*q)/(b*c - a*d)]]        
\end{lstlisting}
       对于如下输入
\begin{lstlisting}
> A:=array( [[a,c],
             [b,d]] );
> u:=array([p,q]);
> linsolve(A,u);
\end{lstlisting}
      \Maple{} 给出如下回答。
      \begin{equation*}
        \bigl[-\frac{-d\,p+q\,c}{-b\,c+a\,d},
          \frac{-b\,p+a\,q}{-b\,c+a\,d}\bigr]
      \end{equation*}
    \end{answer}
\end{exercises}
\index{computer algebra systems|)}
\endinput


